\section{Current Regulatory and Governance Context}

\subsection{Interpretive rule}

This section is current as of 5 August 2026 and uses primary official sources. It maps evidence
features to governance questions; it does not convert a metric into a compliance conclusion.
Applicability depends on the institution, role, jurisdiction, product, use, and effective dates.
Legal and compliance owners must review the facts of a deployment.

\subsection{Framework-by-framework mapping}

The following expanded layout replaces the former compact table so that applicability boundaries
and non-claims remain readable and accessible.

\subsubsection{Federal Reserve model risk}

\textbf{Current state and applicability.} SR 26-2 was issued on 17 April 2026 and supersedes and replaces
SR 11-7 and SR 21-8. The guidance is most relevant to Federal Reserve-supervised banking
organizations with more than \$30 billion in total consolidated assets; use-specific materiality
and institutional scope still matter \citep{sr26_2}.

\textbf{Appropriate evidence use.} Provenance, validation scope, outcome
monitoring, limitations, change control, and independent challenge can support model-risk work.
This paper does not assert conformance, applicability to every institution, or supervisory
acceptance.

\subsubsection{ECOA and Regulation B}

\textbf{Current state and applicability.} The 2026 final rule, effective 21 July 2026, states that ECOA does not
authorize disparate-impact liability or an effects test \citep{cfpb_regb_2026}. Other statutory,
constitutional, contractual, state, and jurisdictional questions may still require review.

\textbf{Appropriate evidence use.} AIR
is used only as a voluntarily selected internal fairness screen and descriptive monitoring ratio.
It is not presented as a Regulation B requirement, an ECOA disparate-impact test, or a lending-law
verdict.

\subsubsection{Employment-selection guidance}

\textbf{Current state and applicability.} The four-fifths convention comes from US employment-selection
guidance. The EEOC describes it as a rule of thumb, not a legal definition, and explains that it is
not controlling in all circumstances \citep{eeoc_uniform_guidelines}.

\textbf{Appropriate evidence use.} The convention explains
the historical source of the paper's 0.80 line and the highest-rate comparison policy. It does not
transfer an employment rule into a lending obligation.

\subsubsection{EU AI Act}

\textbf{Current state and applicability.} Regulation (EU) 2024/1689 includes data-governance, technical-documentation,
transparency, accuracy, robustness, and human-oversight requirements for relevant systems
\citep{eu_ai_act}. The 2026 AI Omnibus entered into force on 27 July 2026; official Commission
materials state 2 December 2027 for Annex III high-risk rules and 2 August 2028 for
product-integrated high-risk rules \citep{eu_ai_omnibus_2026,eu_ai_act_faq_2026}.

\textbf{Appropriate evidence use.} Evidence lineage,
declared data scope, monitoring, robustness, and human-review records may support a case-specific
control framework. This paper does not determine provider/deployer roles, high-risk classification,
conformity, or CE-marking status.

\subsubsection{FCA Consumer Duty}

\textbf{Current state and applicability.} Current FCA publications and FG22/5 materials continue to emphasize monitoring
and evidence of good customer outcomes; perimeter and role remain firm-specific
\citep{fca_consumer_duty_2026}.

\textbf{Appropriate evidence use.} Group-level outcome monitoring, investigation records, and
governance escalation can be supporting evidence. A generated fixture does not demonstrate good
outcomes for customers or FCA compliance.

\subsection{Control mapping, not certification}

Within those limits, the companion can support a control owner who needs to answer narrower
questions:

\begin{itemize}
  \item Which exact product source, runtime image, configuration, dataset hash, and seed produced
        the evidence?
  \item Were protected-group definitions and comparison policies declared before interpretation?
  \item Were unevaluated, non-informative, small-group, utility, and privacy limitations preserved?
  \item Can another reviewer recompute the files, certificate hashes, and displayed numbers?
  \item Is a named human accountable for deciding whether additional investigation is required?
\end{itemize}

The answer to those questions may be useful in a compliance program, but usefulness is not itself
compliance. The paper avoids labels such as ``compliant,'' ``certified,'' ``approved,'' or
``regulator-ready'' for the characterization outcome.

\subsection{Machine-readable dated overlay}

The companion includes
\path{evidence/v5.0.1/publication/regulatory_mapping_2026-08-05.csv}. Each row records the as-of
date, current instrument, applicability boundary, appropriate evidence use, explicit non-claim,
and primary source URL. The original producer regulatory matrix is retained unchanged for evidence
provenance and is not treated as the current publication interpretation.
